\documentclass[a4paper,11pt]{article}

\usepackage{hyperref}
\hypersetup{colorlinks,breaklinks,urlcolor=blue,linkcolor=blue}

\usepackage{libertine}

\usepackage{longtable}

\usepackage[cm]{fullpage}

\usepackage{titlesec}
\titleformat{\section}{\Large\scshape\raggedright}{}{0em}{}[\titlerule]
\titlespacing{\section}{0pt}{1em}{3pt}

\usepackage{lastpage}

\usepackage{fancyhdr}
\pagestyle{fancy}
\fancyhead{}
\renewcommand{\headrulewidth}{0pt}
\fancyfoot[L]{Павел Турянский}
\fancyfoot[C]{Страница \thepage\,из \pageref*{LastPage}}
\fancyfoot[R]{\today}
\renewcommand{\footrulewidth}{0.4pt}

\begin{document}
  \begin{center}
    \Huge \textsc{Павел Турянский}
  \end{center}

  \hrule
  \vspace{0.5em}

  \begin{tabular}{rl}
    \textsc{Telegram:}  & \href{https://t.me/cmcmA20}{@cmcmA20}                         \\
    \textsc{GitHub:}    & \href{https://github.com/cmcmA20}{cmcmA20}                    \\
    \textsc{Codeberg:}  & \href{https://codeberg.org/cmcmA20}{cmcmA20}                  \\
    \textsc{Эл. почта:} & \href{mailto:jpopksad@protonmail.ch}{jpopksad@protonmail.ch}  \\
    \textsc{GPG:}       & C9B0 86E3 C5E8 10F1 794E  E5A4 A65B 7838 0A49 A905            \\
  \end{tabular}

  \section{Краткое описание}
    Я математик-любитель и разработчик программного обеспечения, живу в России. В настоящее время интересуюсь функциональным
    программированием, теорией типов, построением систем с высокими гарантиями.

  \section{Навыки и достижения}
    \begin{longtable}{rp{0.6\textwidth}}
      \textsc{Естественные языки:} & русский, английский \\

      \textsc{Языки программирования:} & Agda, Idris, Haskell, C, Arend, Julia, GNU Octave, Python, x86 asm, nix  \\
                                      & (в порядке убывания владения)                                             \\

      \textsc{Программирование:} & имею $\sim$5 лет опыта работы с Haskell                                                          \\
                                 & внёс вклад в проект \href{https://github.com/agda/agda}{Agda}                                    \\
                                 & поддерживаю \href{https://github.com/cmcmA20/cubical-mini}{собственную стандартную библиотеку}
                                   для Cubical Agda                                                                                 \\
                                 & увлечён системами эффектов и зависимыми типами                                                   \\
                                 & имею небольшой опыт модификации прошивок                                                         \\

      \textsc{Безопасность:} & умею использовать логическое программирование для решения задач безопасности, таких как авторизация
                               и анализ информационных потоков                                                                      \\
                             & могу сделать нечто похожее на действительно безопасную загрузку с coreboot + HEADS и TPM             \\
                             & имею базовый опыт работы с ПЭМИН                                                                     \\

      \textsc{Сети:} & проектировал и контролировал строительство Ethernet-сетей
                       (до сотни ПК, в основном медь и немного оптики)                                           \\
                     & проектировал и контролировал строительство сетей виртуализации                            \\
                     & спроектировал и разработал простую масштабируемую VoIP-систему на Asterisk + собственная
                       маршрутизация вызовов на LUA                                                              \\

      \textsc{Виртуализация:} & проектировал, строил и обслуживал кластеры виртуализации на VMware ESXi и Proxmox  \\
                              & переводил устаревшие инфраструктуры в частное облако                               \\

      \textsc{ОС:} & эксплуатировал Windows Server (2003 - 2012R2), FreeBSD, Linux (Debian, CentOS) в промышленной эксплуатации  \\
                   & лично использовал QubesOS, Void Linux, OpenBSD, пробовал seL4, KasperskyOS                                  \\

      \textsc{Переводы:} & \href{https://www.youtube.com/playlist?list=PLbgaMIhjbmEnaH\_LTkxLI7FMa2HsnawM\_}
                               {Теория категорий для программистов Бартоша Милевского}, см. русские субтитры  \\
    \end{longtable}

  \section{Интересы}
    \begin{tabular}{rl}
      \textsc{Математика:} & изучаю разные вещи из nLab и 1Lab                                          \\
                           & имею опыт работы с HoTT, втч прикладными аспектами                         \\
                           & мечтаю о мультимодальной HoTT с линейностью, когезией и guarded-рекурсией  \\
                           & изучал топологию и дифгем под руководством Искандера Тайманова             \\
                           & прошёл курс Роберта Грайста
                              \href{https://www.coursera.org/account/accomplishments/records/9Y3N9X2BMR}{Calculus:
                              Single Variable}                                                                      \\
      \textsc{Аппаратное обеспечение:} & вожусь с ThinkPad                                          \\
                                       & всё ещё хочу научиться создавать БПЛА и автономное оружие  \\
    \end{tabular}

  \pagebreak
  \section{Опыт работы}
    \begin{longtable}{r|p{0.72\textwidth}}
      \textsc{мар 2026--наст. вр.} & \large \sffamily \textbf{Независимый исследователь}                                     \\
      \nopagebreak
      (в настоящее время) & разрабатываю новые подходы к механизации уровень-полиморфных категориальных структур,
                            см. \href{https://hott-uf.github.io/2026/abstracts/HoTTUF_2026_paper_10.pdf}{тезис HoTT/UF'26}   \\
      \\&\\
      \textsc{июл 2024--мар 2026} & \large \sffamily \textbf{Старший исследователь-разработчик} в \textbf{Kaspersky Lab}, Россия   \\
      \nopagebreak
      (1 год 9 месяцев) & исследовал применения полиномиальных функторов                                                   \\
                         & исследовал форматы передачи данных и языки спецификации протоколов                                     \\
      \\&\\
      \textsc{сен 2021--июл 2024} & \large \sffamily \textbf{Разработчик программного обеспечения} в \textbf{Kaspersky Lab}, Россия   \\
      \nopagebreak
      (2 года 10 месяцев) & разрабатывал KasperskyOS с использованием Haskell и Agda                                            \\
                           & преподавал курс по зависимым типам                                                        \\
      \\&\\
      \textsc{дек 2019--авг 2021} & \large \sffamily \textbf{Разработчик программного обеспечения} в \textbf{Tinkoff Bank}, Россия     \\
      \nopagebreak
      (1 год 9 месяцев) & поддерживал и развивал backend Tinkoff Travel (написан на Haskell)                       \\
                         & разрабатывал ПО потоковой обработки с использованием Apache Kafka                                    \\
                         & научился использовать системы эффектов для структурирования кода на Haskell, в частности, fused-effects   \\
                         & приходилось использовать устаревшие протоколы для взаимодействия с GDS                                        \\
      \\&\\
      \textsc{окт 2017--июл 2019} & \large \sffamily \textbf{Инженер по информационной безопасности} в \textbf{FSUE NPP Gamma},
                                    Россия                                                                               \\
      \nopagebreak
      (1 год 10 месяцев) & дежурил на «Самара Арене» во время Чемпионата мира по футболу FIFA 2018 в качестве консультанта по технической безопасности
                            \\
                          & занимался автоматизацией на PowerShell, управлением конфигурацией физического парка (с Saltstack)   \\
                          & работал со спектроанализаторами Rohde\&Schwarz и Agilent через Python + NIVISA               \\
                          & проектировал компьютерные сети                                                                   \\
                          & приходилось работать с малоизвестным оборудованием и ПО                                              \\
      \\&\\
      \textsc{янв 2017--авг 2017} & \large \sffamily \textbf{Сетевой инженер} в \textbf{Sibirskie Seti Ltd.}, Россия   \\
      \nopagebreak
      (8 месяцев) & настраивал оборудование следующих производителей: Extreme Networks, Juniper, Cisco, D-Link, TP-Link      \\
                 & писал скрипты сетевого мониторинга для Nagios на bash/Python                                         \\
                 & проводил тестирование оборудования и сетей с помощью Python + scapy                                           \\
                 & взаимодействовал с NOC других провайдеров                                                                  \\
                 & однажды получил L2-петлю в сети управления                                                         \\
      \\&\\
      \textsc{сен 2016--дек 2016} & \large \sffamily \textbf{Инженер-техник} в \textbf{ITVaksa Ltd.}, Россия    \\
      \nopagebreak
      (4 месяца) & всё от физической прокладки кабеля до проектирования сетей/систем для малого бизнеса                              \\
    \end{longtable}

  \section{Образование}
    \begin{longtable}{r|p{0.72\textwidth}}
      \textsc{сен 2014--май 2018} & \large \sffamily \textbf{Бакалавр информатики} в \textbf{SibSUTI},
                                    Новосибирск, Россия                                               \\
      \nopagebreak
      (4 года) &                                                                                     \\
      \\&\\
      \iffalse
      \textsc{сен 2006--май 2010} & \large \sffamily в \textbf{Лицее № 1}, Братск, Россия             \\
      \nopagebreak
      (4 года) & физико-математическая специализация                                                     \\
      \fi
    \end{longtable}

\end{document}
